本节给出程序运行结果和最终破解得到的明文。程序首先对目标密文进行分组，并检测原始密文是否能够通过服务器的 PKCS\#7 填充检查。运行结果显示，原始密文 padding 正确，说明解密服务器可以正常访问，Padding Oracle 查询函数可用。

\subsection{Oracle 可用性测试}

程序运行后首先输出目标密文的分组结果：

\begin{lstlisting}[breaklines=true, basicstyle=\small\ttfamily]
IV = 46307250616464316e674f7261636c33
C1 = 9ae0735429869542efc40dcdc3f4c170
C2 = 649463f719c5ddf4ce8c6d1ef0e5d41a
C3 = c5d137629e3fe1340cfaad7e21d65d14
\end{lstlisting}

随后程序向解密服务器提交原始密文进行测试，服务器返回填充正确：

\begin{lstlisting}[breaklines=true, basicstyle=\small\ttfamily]
正在测试 oracle...
原始密文 padding 正确，oracle 可用
\end{lstlisting}

这说明本实验所使用的解密接口能够正常工作，后续可以通过构造密文并观察服务器返回状态进行 Padding Oracle 攻击。

\subsection{明文分组恢复结果}

程序依次对 $C_1$、$C_2$ 和 $C_3$ 三个密文分组进行攻击。每个分组均从最后一个字节开始向前恢复，逐字节得到对应的中间状态和明文字节。

第 1 个明文分组恢复结果如下：

\begin{lstlisting}[breaklines=true, basicstyle=\small\ttfamily]
第 1 个明文分组 hex：
666c61677b3535356d73646236617864

第 1 个明文分组 text：
flag{555msdb6axd
\end{lstlisting}

第 2 个明文分组恢复结果如下：

\begin{lstlisting}[breaklines=true, basicstyle=\small\ttfamily]
第 2 个明文分组 hex：
3962683065716c76683167666378326b

第 2 个明文分组 text：
9bh0eqlvh1gfcx2k
\end{lstlisting}

第 3 个明文分组恢复结果如下：

\begin{lstlisting}[breaklines=true, basicstyle=\small\ttfamily]
第 3 个明文分组 hex：
6379723177733636367d060606060606

第 3 个明文分组 text：
cyr1ws666}
\end{lstlisting}

其中，第 3 个明文分组末尾的：

\begin{equation}
    06\ 06\ 06\ 06\ 06\ 06
\end{equation}

为 PKCS\#7 填充字节，表示最后填充了 6 个字节。因此，真实明文不包含这 6 个填充字节。

\subsection{完整明文恢复结果}

将三个明文分组拼接后，可以得到带有 PKCS\#7 填充的完整明文，其十六进制表示为：

\begin{lstlisting}[breaklines=true, basicstyle=\small\ttfamily]
666c61677b3535356d736462366178643962683065716c76683167666378326b6379723177733636367d060606060606
\end{lstlisting}

去除末尾 6 个字节的 PKCS\#7 填充后，得到真实明文的十六进制表示为：

\begin{lstlisting}[breaklines=true, basicstyle=\small\ttfamily]
666c61677b3535356d736462366178643962683065716c76683167666378326b6379723177733636367d
\end{lstlisting}

将其按 UTF-8 字符串解码，最终得到 5 号目标密文对应的明文为：

\begin{center}
    \fbox{
        \texttt{flag\{555msdb6axd9bh0eqlvh1gfcx2kcyr1ws666\}}
    }
\end{center}

\subsection{运行过程截图说明}

为证明攻击过程确实通过逐字节恢复完成，实验报告中应加入程序运行截图。建议至少插入以下两类截图：

\begin{itemize}
    \item 第一类：Oracle 可用性测试截图，显示目标密文分组以及“原始密文 padding 正确，oracle 可用”；
    \item 第二类：攻击过程截图，显示程序逐字节输出 \verb|fake|、\verb|intermediate| 和 \verb|plain| 的恢复过程；
    \item 第三类：最终结果截图，显示带 padding 的明文、去除 padding 后的明文以及最终明文。
\end{itemize}


